\documentclass[11pt]{article}
\usepackage[T1]{fontenc}
\usepackage{lmodern,amsmath,amssymb,amsthm,booktabs,longtable}
\usepackage[margin=1in]{geometry}
\usepackage[hidelinks]{hyperref}
\usepackage{microtype}
\hypersetup{pdftitle={Explicit Pressure and Local Contraction Bounds for a Forced Navier-Stokes Blow-Up Construction},pdfauthor={Daniel Fredriksen},pdfsubject={Conditional explicit estimates for a leading local analytic profile},pdfkeywords={Navier-Stokes, explicit bounds, contraction mapping}}
\newtheorem{theorem}{Theorem}
\newtheorem{lemma}[theorem]{Lemma}
\newtheorem{proposition}[theorem]{Proposition}
\newcommand{\norm}[1]{\left\lVert#1\right\rVert_\rho}
\newcommand{\Av}{\mathcal A}
\newcommand{\DY}{\mathcal D}
\newcommand{\J}{\mathcal J}
\newcommand{\Int}{\mathcal I}
\title{Explicit Pressure and Local Contraction Bounds\\for a Forced Navier--Stokes Blow-Up Construction}
\author{Daniel Fredriksen\\\small Quantyra}
\date{September 8, 2026\\\small Research note, version 0.1.0}
\begin{document}
\maketitle
\begin{abstract}
We collect explicit sufficient bounds for two local components of the
September 2026 forced Navier--Stokes construction attributed to OpenAI.
An integrable envelope for its unbumped outer schedule gives pressure
majorants $16\overline P^2$ and $128\overline P^2$ on a fixed complex rectangle.
For the leading axis equations, conservative coefficient estimates give a
finite, noncircular rule selecting the contraction parameter and then the
amplitude. The resulting fixed-point map has displacement and Lipschitz
constant at most $1/16$ on a unit ball, and its real normalized angular
profile is at least $1/6$ on the stated interval. These are conditional
quantitative refinements of the source's existing local argument. They do
not select globally admissible outer data, certify the completed velocity
field, or provide a computational complexity result.
\end{abstract}

\section{Contribution, provenance, and scope}
The external manuscript \cite{openai} already introduces the pressure datum,
weighted analytic space, radial inverse estimates, and local contraction.
Our purpose is to expose sufficient constants and their order of dependence.
In particular, this note does not claim a new blow-up theorem or a new
fixed-point method. No priority claim is made for the general techniques.
The source-specific estimates are candidates for a quantitative follow-up;
a focused literature screen is not a certification of novelty.

The derivations were developed in the research records listed in the
repository's provenance manifest. This manuscript consolidates the pressure,
coefficient, and contraction records from one fixed commit. Its mathematical
arguments are informal. Existing AI-assisted reviews concern those records,
not independent human peer review or a machine-checked proof of this note.
The source manuscript is used as an external mathematical input, not
independently certified here. Appendix labels below refer to its 166-page
version retrieved on September 8, 2026; its checksum is recorded separately.

\section{A pressure envelope on a fixed complex domain}
Let $P_*>0$, and let $c:\mathbb R\to(0,\infty)$ and
$\theta:\mathbb R\to[0,1]$ be measurable. Assume the envelope
\begin{equation}\label{env}
c(y)\le P_*\begin{cases}
e^{y/10},&y\le 1,\\
e^{1/10}e^{-(y-1)/2},&y\ge1.
\end{cases}
\end{equation}
These assumptions isolate exactly the radial estimate needed here.
For the source's unbumped schedule, $c'/c=l-1/2$: the inner segment has
$l=3/5$, the first unit interval has $l\le3/5$, and subsequent segments
have $l\le0$ at the central angular parameter. Integrating those inequalities
gives \eqref{env}. This statement concerns the reference schedule in the
datum (A.21), not arbitrary later edited profiles or the complete pressure.

Define
\begin{equation}\label{pressure}
\Pi_0(z)=-\frac12\int_{\mathbb R}c(y)^2
\exp\bigl(-2\theta(y)\operatorname{Log}(1+z^2)\bigr)\,dy,
\quad
\Omega=\{z:|\operatorname{Re}z|<5/4,\ |\operatorname{Im}z|<1/4\}.
\end{equation}
The logarithm is the principal branch.

\begin{theorem}\label{thmpressure}
Under \eqref{env}, $\Pi_0$ is holomorphic on $\Omega$. For every rational
$\overline P\ge P_*$,
\[
|\Pi_0(z)|\le16\overline P^2=:P_0,
\qquad |\Pi_0'(z)|\le128\overline P^2=:P_1
\quad(z\in\Omega).
\]
\end{theorem}
\begin{proof}
The three intervals $(-\infty,0)$, $[0,1]$, and $[1,\infty)$ contribute
at most $5P_*^2$, $e^{1/5}P_*^2$, and $e^{1/5}P_*^2$, respectively, to
$\int c^2$. Since $e^{1/5}\le\sum_{k\ge0}(1/5)^k=5/4$,
\[
\int_{\mathbb R}c^2\le(5+2e^{1/5})P_*^2<8P_*^2.
\]
On $\Omega$, $\operatorname{Re}(1+z^2)>15/16$ and $|z|<3/2$.
Consequently $f=(1+z^2)^{-1}$ satisfies $|f|<2$ and $|f'/f|<4$.
For real $0\le\theta\le1$, the integrand's power has magnitude at most
$4$ and its $z$ derivative has magnitude at most $32$.
These bounds give common integrable majorants. Holomorphic integration and
differentiation under the integral are therefore valid, and multiplication
by the factor $1/2$ proves the two bounds.
\end{proof}

For $R\ge1$, the pressure contribution outside $[-R,R]$ is at most
\begin{equation}\label{tail}
10\overline P^2e^{-R/5}+3\overline P^2e^{-(R-1)}.
\end{equation}
The corresponding derivative tail is at most eight times this expression.
This follows by integrating the two tails in \eqref{env} and using the
same power bounds. For a finite schedule with computable parameters and
effective compact-interval moduli, certified quadrature and \eqref{tail}
give finite-accuracy pressure evaluation. Measurability and the envelope
alone do not imply computability. No quadrature implementation or cost
bound is asserted here.

\section{Coefficient domain and explicit local data}
Fix rational parameters
\[
0<h\le1/100,\quad 0<j\le1/20,\quad 0<\sigma\le1,\quad
0\le e\le1/1000,
\]
and choose rational $0<r\le\min(1/1000,\sigma^2/4096)$.
Write $I=[-1-e,1+e]$, let $\Omega_r$ be its open complex $r$-neighborhood,
and put $\rho=r/4$. Define
\begin{align*}
A&=1/2+h,& D&=1/2-h,& d&=1-\eta^2,& L&=1-2h\eta^2,\\
U_*&=4\eta+j,& H_*&=D\eta+dU_*,& Q_*&=H_*^2+\sigma^2,\\
\chi&=H_*^2/Q_*,& \zeta&=-LH_*/Q_*.
\end{align*}
\[
W_*=1-dU_*'-2D\eta U_*,\qquad
Z_*=-A(1-2\eta U_*)U_*-H_*U_*'-d\Pi_0'+4A\eta\Pi_0.
\]
These are the coefficients in the source's Appendix B. The source's
global outer admissibility and additional separation condition (B.2) are
not established by the parameter ranges above. Those remain additional
obligations when applying the local construction to the complete source.

\begin{lemma}\label{domains}
On $\Omega_r$ one has
\[
|H_*|<9,\quad |H_*'|<18,\quad 1/2<|L|<2,\quad |Q_*|\ge\sigma^2/2.
\]
Furthermore,
\[
|\chi|\le162/\sigma^2,\quad |\zeta|\le36/\sigma^2,\quad
|Z_*/L|\le\overline Z:=192+928\overline P^2.
\]
\end{lemma}
\begin{proof}
The tube lies in $|z|<101/100$. Triangle bounds for the displayed
polynomials give the first three inequalities and
$|(H_*^2)'|<324$. For each $z$ in the tube there is real $x\in I$
with $|z-x|<r$. Since $Q_*(x)\ge\sigma^2$, integration along the
segment gives $|Q_*(z)|\ge\sigma^2-324r>\sigma^2/2$.
This proves the bounds for $\chi$ and $\zeta$.
For $Z_*$, use $|U_*|<5$, $|d|<3$, $|1-2zU_*|<12$,
$|H_*U_*'|<36$, $A<1$, and $|4Az|<5$.
It follows that $|Z_*/L|\le192+10P_0+6P_1=\overline Z$.
The entire tube lies inside $\Omega$, so Theorem~\ref{thmpressure} applies.
\end{proof}

For $\Lambda\ge1$ set
\[
\phi_*(z)=\exp\left(\Lambda\int_0^z\zeta(w)\,dw\right),\qquad
g=\phi_*/C.
\]
The tube is convex, the path length is less than two, and hence
$|\phi_*|\le e^{72\Lambda/\sigma^2}$ there. In particular,
\begin{equation}\label{Cchoice}
C\ge2^{\lceil144\Lambda/\sigma^2\rceil}
\quad\Longrightarrow\quad |g|\le1\text{ on }\Omega_r,
\end{equation}
because $\log2>1/2$. A bound on this complex domain, rather than only
on its real interval, is needed to control coefficient derivatives.

\section{The weighted space and operator constants}
We use the source's space with the following normalization. Write
$F(Y,\eta)=\sum_{\alpha\ge0}F_\alpha(\eta)Y^\alpha$ and put
\[
a_{\alpha\beta}=
\frac{20^{-\alpha}\rho^{-\beta}\beta!\binom{\alpha+\beta}{\beta}}
{(\alpha+1)^2(\beta+1)^2},\qquad
\norm{F}=\sup_{\alpha,\beta\ge0,\ \eta\in I}
\frac{|\partial_\eta^\beta F_\alpha(\eta)|}{a_{\alpha\beta}}.
\]
The space $B_\rho$ consists of smooth coefficient sequences with finite
norm and is complete, by uniform convergence at every coefficient and
derivative order. On $B_\rho^2$ use the maximum norm. Set
\[
\Int F=\int_0^Y F(s,\eta)\,ds,\quad \Av F=Y^{-1}\Int F,
\quad \DY=Y\partial_Y,
\quad (\J_\nu F)_{\alpha+1}=\frac{F_\alpha}{(\alpha+1)(\alpha+\nu)},
\quad (\J_\nu F)_0=0.
\]
Here $\nu=1,2$; averaging is defined at zero by its coefficient extension.

The coefficient convolution and weight shifts in Lemma B.1 give
\begin{equation}\label{ops}
\begin{gathered}
a_{\rm alg}=256,\quad \|\J_\nu\|,\|\Int\|,\|Y\cdot\|\le80,
\quad \|\partial_\eta\Int\|\le80/\rho,\quad \|\Av\|\le1,\\
\norm{\J_\nu[(\partial_\eta F)(\DY G)]}
\le C_{\rm mix}\norm{F}\norm{G},\qquad C_{\rm mix}=20480/\rho.
\end{gathered}
\end{equation}
For completeness, the convolution constant in each index is bounded by
$8\sum_{i\ge1}i^{-2}<16$, giving $256$ for multiplication.
The weight ratios satisfy
\[
\frac{a_{\alpha\beta}}{a_{\alpha+1,\beta}}\le80,
\qquad \frac{a_{i,k+1}}{a_{i+1,k}}\le\frac{80}{\rho}(i+1).
\]
For a mixed term at output radial degree $N=\alpha+1$, assign the
raised degree to the factor carrying $\partial_\eta$. The residual factor
after radial inversion is
$(i+1)(\alpha-i)/[N(\alpha+\nu)]\le1$. If the radial derivative is
omitted, replace $\alpha-i$ by one. Without the parameter derivative,
the first weight ratio gives $80(\alpha-i)/[N(\alpha+\nu)]\le80$;
omitting both is easier. Thus the same $C_{\rm mix}$ covers all these
cases, including $\alpha=0$. Averaging costs no factor. Extra
undifferentiated coefficients are handled by further convolution, each
at cost $256$; they are not pulled through radial inversion.
This does not assert boundedness of either unintegrated derivative alone.

For an $\eta$-only function bounded by $M$ on $\Omega_r$, Cauchy estimates
on disks of radius $r/2$ give
\[
\norm{a}\le M\sup_{\beta\ge0}\frac{(\beta+1)^2}{2^\beta}\le3M.
\]
Consequently multiplication by $\chi$ has norm at most
$M_\chi=124416/\sigma^2$. If $T_\chi=\J_2\chi/2$, where multiplication
precedes integration, its radial degree increases at every application.
The sharper bound on inputs with no degree below $b$ is
$\|\J_2F\|\le80\|F\|/[(b+1)(b+2)]$, so
\[
\|T_\chi^k\|\le\frac{(40M_\chi)^k}{k!(k+1)!},\qquad
\|(1+T_\chi)^{-1}\|\le e^{40M_\chi}\le V:=2^{\lceil80M_\chi\rceil}.
\]
The inverse follows from the absolutely convergent alternating operator
series. It does not require $\|T_\chi\|<1$.

\section{Integrated remainders and a contraction rule}
Set
\begin{equation}\label{constants}
\begin{gathered}
K=\max(V,120\overline Z)+1,\qquad B_c=1000/\sigma^2,\\
T=a_{\rm alg}B_c C_{\rm mix},\qquad
P_c=57600a_{\rm alg}^4B_c(2+1/\rho),\\
N=T(6K+11K^2)+P_cK^2,\qquad
L_R=T(6+22K)+2P_cK.
\end{gathered}
\end{equation}
All these constants are fixed before $\Lambda$ and $C$.
The center is
\[
\Phi_0=f_0(Y\chi),\quad u_0=-YZ_*/(2L),\qquad
f_0(z)=\sum_{\alpha\ge0}\frac{(-z/2)^\alpha}{\alpha!(\alpha+1)!}.
\]
Its component norms are at most $V$ and $120\overline Z$, respectively,
so the closed unit ball about it has both component norms at most $K$.

Write $\varepsilon=\Lambda^{-1}$ and define
\[
B=-2D\eta\Av u-d\partial_\eta\Av u,\quad
W=W_*+\varepsilon B,\quad U=U_*+\varepsilon u,\quad
H_c=H_*+\varepsilon du,\quad p=\Int(g^2\Phi^2).
\]
The source remainders, with their full derivative dependence, are
\begin{align}
LR_1={}&[W+h(1-2\eta U)+du\zeta]\Phi+W\DY\Phi+H_c\partial_\eta\Phi,
\label{r1}\\
LR_2={}&[A(1-4\eta U_*)+dU_*']u-2A\eta\varepsilon u^2
+W\DY u+H_*\partial_\eta u+\varepsilon du\partial_\eta u\notag\\
&\hspace{8mm}-4A\eta p+d\partial_\eta p-2\eta\DY p.\label{r2}
\end{align}
Expanding these expressions gives three linear and seven quadratic
nonpressure terms in $R_1$, and three linear, four quadratic, and three
pressure terms in $R_2$. Appendix~\ref{terms} lists them individually.

Every $\eta$-only coefficient in that list, including its division by $L$,
has norm at most $B_c$. For example $|W_*|<19$ on the tube, and
$\norm{d\zeta/L}\le648/\sigma^2$. The axial linear coefficient before
division by $L$ has magnitude below $34$; Cauchy conversion bounds its
norm after division by $204<210$. The other coefficients obey smaller
bounds by Lemma~\ref{domains}. After applying the indicated $\J_\nu$,
\eqref{ops} therefore gives $TK$ per nonpressure linear term and $TK^2$
per quadratic term. The powers of $\varepsilon$ are at most one.

By \eqref{Cchoice}, $\norm{g}\le3$. For one fixed $g$ per contraction,
four-factor multiplication gives
\[
\norm{g^2\Phi^2}\le9a_{\rm alg}^3K^2,\qquad
\norm{g^2(\Phi^2-\Psi^2)}\le18a_{\rm alg}^3K\norm{\Phi-\Psi}.
\]
The operators generating $p$, $\partial_\eta p$, and $\DY p$ from
$g^2\Phi^2$ are $\Int$, $\partial_\eta\Int$, and multiplication by $Y$.
Applying these first, then coefficient multiplication and the outer
$\J_\nu$, gives the total pressure bound
$80\cdot80\cdot9a_{\rm alg}^4B_c(2+1/\rho)K^2=P_cK^2$.
Replacing the two unknown factors one at a time gives its difference
bound $2P_cK$ times the input difference.
Summing both component bounds proves that $N$ and $L_R$ in
\eqref{constants} bound the norm and Lipschitz constant of
$(\J_2R_1,\J_1R_2)$ on the unit ball.

\begin{theorem}\label{contraction}
For the local data and pressure above, choose an integer
\[
\Lambda\ge\max\{1,\lceil8V(N+L_R+1)\rceil\},
\]
and then choose $C$ satisfying \eqref{Cchoice}.
The map
\[
\mathcal F(\Phi,u)=\left(
\Phi_0+\frac{(1+T_\chi)^{-1}\J_2R_1}{2\Lambda},\quad
u_0+\frac{\J_1R_2}{2\Lambda}\right)
\]
maps the closed unit ball about $(\Phi_0,u_0)$ into itself and has
Lipschitz constant at most $1/16$. Its unique fixed point in that ball
has distance at most $1/16$ from the center. On $0\le Y\le41/10$,
$\eta\in I$, its real angular component satisfies $\Phi\ge1/6$.
\end{theorem}
\begin{proof}
The map's displacement and Lipschitz constant are bounded by
$VN/(2\Lambda)$ and $VL_R/(2\Lambda)$, respectively, each at most
$1/16$. Banach's theorem applies on the complete closed ball.
Although $g$ depends on the later choices, its common norm bound follows
from the same larger complex domain for every allowed pair $(\Lambda,C)$.
Thus selecting $\Lambda$ using these constants is noncircular.

The radial-zero conditions $\Phi(0,\eta)=1$ and $u(0,\eta)=0$ are preserved:
both radial inverses have zero constant coefficient, and the angular
inverse preserves this subspace term by term. Real data and iteration from
the real center give real coefficients.
For $\beta=0$, the weights imply
$|F(Y,\eta)|\le\norm{F}/(1-|Y|/20)$, whose factor is below $4/3$
on the stated interval. Also $0\le\chi\le1$ on $I$.
The alternating series for $f_0$ gives, with $t=z/2$,
\[
f_0(z)\ge1-t/2+t^2/12-t^3/144\ge305719/1152000>1/4
\quad(0\le z\le41/10).
\]
The cubic is decreasing on $[0,41/20]$; the remaining alternating
tail starts positively with decreasing magnitudes. Therefore
$\Phi\ge1/4-(4/3)(1/16)=1/6$.
\end{proof}

The coefficient binomial sum also gives convergence on $|Y|\le5$ with
complex $\eta$ displacement less than $\rho/4$ from the real interval.
Indeed the geometric factors have sum at most $1/4+1/4<1$ before
the polynomial weight denominators are used. In the original radial
variable this establishes a leading local profile on
$0\le X\le41/(10\Lambda)$. It does not establish the source's later
continuation or endpoint shear inequalities.

\section{Effectivity, verification, and remaining obligations}
Exact Picard iterates starting at the center have first displacement at
most $1/16$ and successive increments bounded by powers of $1/16$.
After $m$ iterations the norm error is at most
$(1/16)^{m+1}/(1-1/16)$. This is a bound for exact analytic iterates.
Finite coefficient truncations need not converge in the full weighted
supremum norm: their tails can continue to attain that norm.
On smaller compact analytic domains, the coefficient envelope instead
gives effective value and fixed-derivative tails. With effective pressure
and parameter data this supports finite-accuracy local evaluation in
principle; an implementation must account for arithmetic error, coefficient
access, and both truncation and iteration errors.

The constants are deliberately loose. Writing a power of two as an
expression is distinct from writing its full binary expansion, and the
number of exact Picard steps is distinct from the work needed to evaluate
each step. Neither a polynomial bit-time bound nor an implemented certified
local solver is provided.

The companion exact-rational script checks elementary numerical
inequalities used in these estimates. It is not a formal proof of the
analytic argument. The archived reviews separately address source
alignment and the conditional claims; the consolidated manuscript has
not received independent human peer review.

To obtain a certified completed field one would still need globally
admissible outer parameters, the source's separation and continuation
conditions, higher-order corrections, and normalized cutoff-tail bounds.
Theorems~\ref{thmpressure} and \ref{contraction} do not discharge those
obligations. They also do not construct logic, memory, a readable fluid
device, or a SAT algorithm. The intended contribution is the explicit
local certificate and its auditable dependence on the stated inputs.

\appendix
\section{Expanded remainder inventory}\label{terms}
Each coefficient below includes division by $L$. A dash in the last
column means a multiplier of one. This inventory fixes the term counts
used in \eqref{constants}.
\begin{center}\small
\begin{tabular}{lll}
\toprule Coefficient in $R_1$ & Unknown factors & Multiplier\\\midrule
$(W_*+h(1-2\eta U_*))/L$ & $\Phi$ & --\\
$d\zeta/L$ & $u\Phi$ & --\\
$W_*/L$ & $\DY\Phi$ & --\\
$H_*/L$ & $\partial_\eta\Phi$ & --\\
$-2D\eta/L$ & $(\Av u)\Phi$ & $\varepsilon$\\
$-d/L$ & $(\partial_\eta\Av u)\Phi$ & $\varepsilon$\\
$-2h\eta/L$ & $u\Phi$ & $\varepsilon$\\
$-2D\eta/L$ & $(\Av u)\DY\Phi$ & $\varepsilon$\\
$-d/L$ & $(\partial_\eta\Av u)\DY\Phi$ & $\varepsilon$\\
$d/L$ & $u\partial_\eta\Phi$ & $\varepsilon$\\\bottomrule
\end{tabular}
\end{center}
\begin{center}\small
\begin{tabular}{lll}
\toprule Coefficient in $R_2$ & Unknown expression & Multiplier\\\midrule
$[A(1-4\eta U_*)+dU_*']/L$ & $u$ & --\\
$W_*/L$ & $\DY u$ & --\\
$H_*/L$ & $\partial_\eta u$ & --\\
$-2A\eta/L$ & $u^2$ & $\varepsilon$\\
$-2D\eta/L$ & $(\Av u)\DY u$ & $\varepsilon$\\
$-d/L$ & $(\partial_\eta\Av u)\DY u$ & $\varepsilon$\\
$d/L$ & $u\partial_\eta u$ & $\varepsilon$\\
$-4A\eta/L$ & $p$ & --\\
$d/L$ & $\partial_\eta p$ & --\\
$-2\eta/L$ & $\DY p$ & --\\\bottomrule
\end{tabular}
\end{center}

\section*{Author and assistance disclosure}
Daniel Fredriksen directed the research and publication preparation.
OpenAI Codex assisted with derivations, internal reviews, consolidation,
and repository preparation. AI-assisted checks are not human peer review;
author responsibility and source attribution remain explicit.

\begin{thebibliography}{9}
\bibitem{openai} OpenAI. \emph{Finite time blowup for Navier--Stokes}.
September 2026 manuscript, 166 pages. Appendix A.2--A.4, especially
(A.21), and Appendix B.1--B.3, especially (B.4)--(B.16).
\href{https://cdn.openai.com/pdf/32d9f210-8b73-45e0-91bc-82a30aef8a9a/navier-stokes.pdf}{Source PDF}.
Retrieved September 8, 2026. Exact downloaded-file checksum and derivation
provenance appear in the companion repository.
\bibitem{companion} Daniel Fredriksen.
\emph{Navier--Stokes effective bounds: companion sources and verification}.
\url{https://github.com/Quantyra/navier-stokes-effective-bounds}.
\end{thebibliography}
\end{document}
